\section{Threats to Validity}\label{sec:threats}
\subsection{Simulator validity and policy comparability}
The primary findings are conditional on a project-local, deterministic,
fixed-step simulator. Native partial-prefill execution, allocation refusal,
KV preemption and recomputation are not reproduced merely by a policy's
name or unit-test success. Historical policies can maintain shadow progress
and KV tables that differ from shared execution state. Some resource limits
are policy-dependent; demand above nominal KV capacity is possible. The
simulator is therefore a comparison of these implemented abstractions,
not a controlled physical comparison isolating only production scheduler
selection rules. Differences in abstraction can influence rankings.

Later source audits reproduced discrepancies in final-prefill sampling,
output-prefix retention under recomputation and admission accounting.
Native scheduling counters are not execution acknowledgments. Intermediate
prefill can produce no sample, while final prefill can produce the first
output token. Restarting generated progress changes work and first-token
accounting. No correction factor translates historical simulated timing
to physical timing; no numerical ranking effect was measured by those audits.

\subsection{Physical evidence boundary}\label{sec:real-system}\label{sec:real-system-objective}\label{sec:real-system-design}\label{sec:real-system-result}
The legacy RQ6 artifact records 240 physical executions of two configured
policy paths across three sources and 40 windows per source. A retrospective
hardware audit established 10~GB A100 MIG allocations, with 9.5~GiB visible
to CUDA, rather than full 80~GB devices. Parent-card names were not
allocated-device identity evidence. Individual allocated UUIDs and exact
compute-slice profiles were not preserved.

More seriously, request class metadata was not transported to the legacy
SLAI plugin; missing classes selected fallback behavior. Its deadline clock
used a synthetic iteration-to-time conversion. Client worker limits and
prompt construction could delay actual arrivals, which were not fully
instrumented, and adapter plans could differ from native scheduled work.
The saved numerical observations remain historical evidence of that runner.
They do not establish faithful SLAI-versus-native ranking agreement or
reversal non-reproduction. The mismatch cannot be attributed uniquely to
simulator error. RQ6 remains unresolved at the intended fidelity level.

A later full-A100 engineering smoke established bounded native/checked-SLAI
serving on short synthetic fixtures. It is separate from this campaign and
does not validate realistic-trace ranking portability. Subsequent CPU ledger
and native-boundary conformance work uses synthetic execution replies;
it is not GPU timing, numerical-output or performance evidence. The P2B
control hook rejects selected waiting work preceding selected resident
work in one batch and has no atomic native-state rollback on late allocation
failure. General partial-prefill serving, KV pressure/preemption,
recomputation and distributed TP2/DP2 remain separate gates. Nothing in
this manuscript establishes cross-architecture or parallelism portability.

\subsection{Statistical and construct validity}
The 18,720 executions comprise 9,360 deterministic pairs, not stochastic
replications. Window bootstrap intervals characterize variation within the
sampled trace corpus under a window-exchangeability assumption. Production
traces can retain temporal or session dependence across windows. Independent
publication of traces is not independent sampling of providers; BurstGPT and
Azure share infrastructure lineage. Provider and domain effects are not
causally isolated. The temporal definitions differ by source and should
not be pooled into a common domain-shift statistic.

Missing completion-conditioned observations are excluded, not imputed.
An estimable aggregate does not imply that every policy/window observation
exists. In the primary panel, 48 unique configurations have zero completions;
the preserved calibration manifest has none. Earlier prose reporting 301 calibration zero-completion cells is disputed by that file. Survivorship
conditioning, deterministic tie breaks and synthetic deadlines affect the
meaning of a preferred policy. Practical-margin and recovery thresholds
are reporting conventions, not universal systems standards. The raw-CI
reversal taxonomy and the additional within-family BH-FDR flag answer
different questions; neither is a family-wise error guarantee. No global
multiplicity claim across all analyses is made.

The cross-metric and SLO extensions were conducted after primary results
were available. Their frozen extension configurations support reproducibility,
not retrospective preregistration. The descriptor regressions are sparse
and descriptive. The finite-reference recovery curves do not prove that
40 windows suffice for a new source or that a full-reference ordering is a
population truth. Alternative requests-per-window sizes, held-out sources
and a prospective transfer predictor are not evaluated.

\subsection{Coverage, provenance and reproducibility}
Source token counts may represent different context semantics and are not
necessarily counts under a common tokenizer. Long prompts in the Azure and
Bailian windows are outside the later 256-token engineering contract. No
retokenization or context-eligibility audit can be inferred from the short
smoke. The 120 selected windows are a finite corpus, not a representative
sample of every production LLM service.

Historical source, campaign and analysis files are preserved with hashes.
M1 independently reconciles execution counts, numerical pairs, primary
Kendall point estimates, stored reversal classifications and cross-metric
point estimates. It does not regenerate every bootstrap interval, establish
fresh hardware reproducibility, or retroactively validate source licensing
and public release accessibility. Bibliography metadata and the prior
manuscript overlap boundary require final author review. Unresolved
checks are listed explicitly in the accompanying claim-to-evidence register.
